\begin{thebibliography}{9}

\bibitem{annotated}
《注解版 Rudin》（\emph{The Annotated Rudin}）。\cite{rudin} 的逐字
版本，附评注、插图与各章附录，附录取材于
\cite{taoI,taoII,yupin,ecnu,ustc}。\S\ref{sec:int} 中用到的振幅引理
在第 2 章与第 4 章的附录中给出证明。

\bibitem{rudin}
W.~Rudin, \emph{Principles of Mathematical Analysis}, 3rd ed.,
McGraw--Hill, 1976.

\bibitem{taoI}
T.~Tao, \emph{Analysis I}, 3rd ed., Hindustan Book Agency, 2014.

\bibitem{taoII}
T.~Tao, \emph{Analysis II}, 3rd ed., Hindustan Book Agency, 2014.

\bibitem{yupin}
于品，《数学分析讲义》，清华大学丘成桐数学科学中心。

\bibitem{ecnu}
华东师范大学数学系，《数学分析》第 5 版，高等教育出版社，2019。

\bibitem{ustc}
程艺，《数学分析讲义》，中国科学技术大学少年班，2005。

\end{thebibliography}
